% ====================================================================
% ABSTRAK (BAHASA INDONESIA)
% ====================================================================
% File ini berisi abstrak dalam bahasa Indonesia (200-300 kata).
% EDIT konten sesuai dengan penelitian Anda.
% ====================================================================

\frontmatterchapter{ABSTRAK}

% --- ISI ABSTRAK ---
\noindent Tingginya aktivitas pertukaran pesan pada grup WhatsApp kerap memicu fenomena \textit{information overload}, yang menyulitkan pengguna dalam mengikuti dinamika diskusi dan menyaring informasi penting tanpa membaca riwayat pesan secara manual. Kendati teknik \textit{topic modeling} dapat mengekstraksi topik secara otomatis, penerapannya pada percakapan WhatsApp menghadapi tantangan teks pendek, struktur informal, derau tinggi, dan \textit{code-mixing}. Di samping itu, algoritma pengklasteran HDBSCAN pada \textit{framework} BERTopic standar mengeliminasi hingga 41\%--43\% data sebagai \textit{outlier}, sehingga menghilangkan banyak substansi percakapan penting.

Penelitian ini bertujuan mengembangkan modifikasi arsitektur BERTopic yang memadukan model \textit{embedding} IndoBERTweet, algoritma pengklasteran hierarkis BIRCH untuk menekan proporsi \textit{outlier}, serta skema pembobotan BM25 untuk representasi kata kunci topik. Model dievaluasi terhadap \textit{baseline} BERTopic menggunakan metrik \textit{Topic Coherence} (C-NPMI), \textit{Topic Diversity}, \textit{Embedding Density}, dan \textit{Intra-topic Similarity}, serta diimplementasikan ke dalam aplikasi analisis berbasis web.

Metodologi penelitian menerapkan pendekatan eksperimental dan R\&D menggunakan dataset 91.868 pesan grup WhatsApp Komunitas Buku (2023--2026) yang menghasilkan 56.089 dokumen bersih setelah \textit{Single-Path Preprocessing}. Pemodelan topik memadukan IndoBERTweet (768 dimensi), reduksi UMAP (5 dimensi), pengklasteran BIRCH, ekstraksi kata kunci BM25 dengan penyaringan \textit{stopwords} 9 kategori, serta reduksi topik otomatis. Sistem dibangun dengan arsitektur \textit{client-server} (Vue.js dan FastAPI) yang dilengkapi pemantauan progres \textit{real-time} berbasis SSE, penelusuran konteks pesan (\textit{traceability}), dan privasi data di sisi klien.

% --- KESIMPULAN (Belum diisi sesuai instruksi) ---
% Kesimpulan dan temuan empiris final akan dilengkapi setelah evaluasi menyeluruh.

\vspace{0.3cm}

% --- KATA KUNCI ---
\noindent \textbf{Kata Kunci:} \textit{BERTopic, IndoBERTweet, BIRCH Clustering, BM25, Topic Modeling, Analisis Grup WhatsApp, Pemrosesan Bahasa Alami}

\newpage
